\documentclass[11pt,a4paper]{article}

%%% Place your CSV files in the same folder as this file, and list them below.
\def\CSVfiles{%
	{Example1.csv},%
	{Example2.csv}% ensure the last one has NO comma after
}

%%% PACKAGES ----------------------------------
\usepackage{geometry}
\geometry{landscape,margin=0.2in}

\usepackage{xcolor}
\usepackage{tcolorbox}

\usepackage{graphicx}
\usepackage{pgffor}
\usepackage{tikz}
\usetikzlibrary{shapes.geometric}

\usepackage{xint}


\usepackage{adjustbox}
\usepackage{catchfile}

\usepackage{etoolbox}
\usepackage{ifthen}

\usepackage{calc}
\usepackage{xstring}
\usepackage{xintexpr}

\usepackage[outline]{contour}

%%% COLORS ----------------------------------
\definecolor{s}{RGB}{173, 255, 47}
\definecolor{h}{RGB}{255, 223, 0}  % h = wheat
\definecolor{g}{RGB}{255, 223, 0}  % g = grain = wheat
\definecolor{w}{RGB}{125,125,0}    % w = wood
\definecolor{l}{RGB}{125,125,0}    % l = lumber = wood
\definecolor{b}{RGB}{227, 66, 52}
\definecolor{o}{RGB}{181, 181, 181}
\definecolor{d}{RGB}{255,255,255}
\definecolor{a}{RGB}{0,255,255}

%%% TIKZ STUFF ----------------------------------
\tikzset{hexsty/.style={draw,shape=regular polygon,regular polygon sides=6,minimum size=42mm,
		inner sep=0pt,outer sep=0pt,shape border rotate=90}}
\tikzset{portsty/.style={draw=blue,circle}}

%%% SYMBOLS FOR LEGEND ----------------------------------
\def\RoadSym{\tikz\fill[scale=0.4] (0, 0) rectangle ++(1.0,0.3) ;} 
\def\SettSym{\tikz\filldraw[scale=0.4] (0, 0) circle (0.3) ;} 
\def\CitySym{\begin{tikzpicture}%
	\filldraw[scale=0.4] (0, 0) circle (0.3) ;
	\draw[scale=0.4, line width=1pt] (0, 0) circle (0.5) ;
	\end{tikzpicture}%
}
\def\DevSym{\begin{tikzpicture}%
	\draw[scale=0.4, line width=0.8pt] (0.6, 0.5) circle (0.2) ;
	\draw[scale=0.4, line width=1pt] (0, 0) rectangle (1.2, 1) ;
	\end{tikzpicture}%
}


%%% TILE AND PORT HELPER FUNCS ----------------------------------
\newcommand{\genResWord}[1]{% creates word on tiles
	\ifthenelse{\equal{#1}{w}}{\contour{white}{\normalsize WOOD}}{}%
	\ifthenelse{\equal{#1}{l}}{\contour{white}{\normalsize LUMBER}}{}%
	\ifthenelse{\equal{#1}{h}}{\contour{white}{\normalsize WHEAT}}{}%
	\ifthenelse{\equal{#1}{g}}{\contour{white}{\normalsize GRAIN}}{}%
	\ifthenelse{\equal{#1}{o}}{\contour{white}{\normalsize ORE}}{}%
	\ifthenelse{\equal{#1}{b}}{\contour{white}{\normalsize BRICK}}{}%
	\ifthenelse{\equal{#1}{s}}{\contour{white}{\normalsize SHEEP}}{}%
	\ifthenelse{\equal{#1}{d}}{\contour{white}{\normalsize DESERT}}{}%
}


\newcommand{\pip}{\raisebox{0.7ex}{\contour{white}{\scriptsize$\bullet$}}}  % a pip (dot)
\newcommand{\genPips}[1]{% generate pip (dots) that correspond to dice roll probability
	\ifnum #1>1% only show pips if 2 or more
	\if\relax\detokenize{#1}\relax%
	\else%
	\edef\Npip{\xintthe\xintiexpr6-abs(#1-7)\relax}%
	\edef\Nspace{\xintthe\xintiexpr5-\Npip\relax}%
	\foreach \n in {1,...,\Npip}{\pip}%
	\foreach \n in {1,...,\Nspace}{\phantom{\pip}}%
	\fi%
	\else%
	\phantom{\pip}
	\fi%
}

\newcommand{\genNum}[1]{ % generate number for tile, underlines 6 or 9, white 0 if desert
	\if\relax\detokenize{#1}\relax%
	% arg is empty
	\else%
	\ifnum #1=6%
	\underline{#1}%
	\else%
	\ifnum #1=9%
	\underline{#1}%
	\else%
	\ifnum #1=0% if 0 dont show anything
	\textcolor{white}{0}
	\else%
	#1%
	\fi%
	\fi%
	\fi%
}

%%% TILE AND PORT FUNCTIONS ----------------------------------
\newcommand{\Tile}[3]{% {RESOURCE}{NUMBER}{GRAPHIC}
	\tcbox[colback=#1,colframe=#1,
	boxsep=1pt,left=4pt,right=4pt,top=4pt,bottom=4pt,
	circular arc]{\centering%
		\parbox[c][23mm][c]{23mm}{
			\mbox{\tcbox[colback=white,colframe=white,
				boxsep=1pt,left=2pt,right=2pt,top=2pt,bottom=2pt,
				circular arc]{\parbox[c][8mm][c]{8mm}{\centering\Huge\genNum{#2}
				}}\small{ \rm \genPips{#2}
			}}
			% TODO if you want, put a graphic here
			% #3 
			\centering \genResWord{#1} \centering
	}}%
}


\newcommand{\Port}[2][]{% if optional arg passed, don't put a 2 or 3 (for port trading)
	\ifthenelse{\equal{#1}{}}{%
	\let\TWO2%
    \let\THREE3%
	}{%
	\let\TWO\empty%
    \let\THREE\empty%
	}%
	\ifthenelse{\equal{#2}{w}}{\small\colorbox{#2}{\contour{white}{\textsuperscript{\TWO}WD}}}{}%
	\ifthenelse{\equal{#2}{l}}{\small\colorbox{#2}{\contour{white}{\textsuperscript{\TWO}LM}}}{}% lumber = wood
	\ifthenelse{\equal{#2}{b}}{\small\colorbox{#2}{\contour{white}{\textsuperscript{\TWO}BR}}}{}%
	\ifthenelse{\equal{#2}{s}}{\small\colorbox{#2}{\contour{white}{\textsuperscript{\TWO}SH}}}{}%
	\ifthenelse{\equal{#2}{h}}{\small\colorbox{#2}{\contour{white}{\textsuperscript{\TWO}WH}}}{}%
	\ifthenelse{\equal{#2}{g}}{\small\colorbox{#2}{\contour{white}{\textsuperscript{\TWO}GR}}}{}% grain = wheat
	\ifthenelse{\equal{#2}{o}}{\small\colorbox{#2}{\contour{white}{\textsuperscript{\TWO}OR}}}{}%
	\ifthenelse{\equal{#2}{a}}{\small\colorbox{#2}{\contour{white}{\textsuperscript{\THREE}??}}}{}%
					}



%%% INIT COUNTERS FOR GENERATING TILE COMMANDS ----------------------------------
\newcounter{index}%
\newcounter{column}%
\newcounter{tilenum}%
	

\begin{document}
				
	\newcommand{\ExtraHeight}{\vphantom{\Huge K}} % for legened tables extra spacing
	\newcommand{\PrintCatanBoard}[1][]{\newpage
		\begin{minipage}{.7\textwidth}
			
			\begin{tikzpicture} % the board template
			
			\node[hexsty,anchor=0]   (TileA) at (0,0)       {\TileA};
			\node[hexsty,anchor=150] (TileB) at (TileA.30)  {\TileB};
			\node[hexsty,anchor=150] (TileC) at (TileB.30)  {\TileC};
			
			\node[hexsty,anchor=90]  (TileD) at (TileA.210) {\TileD};
			\node[hexsty,anchor=150] (TileE) at (TileD.30)  {\TileE};
			\node[hexsty,anchor=150] (TileF) at (TileE.30)  {\TileF};
			\node[hexsty,anchor=150] (TileG) at (TileF.30)  {\TileG};
			
			\node[hexsty,anchor=90]  (TileH) at (TileD.210) {\TileH};
			\node[hexsty,anchor=150] (TileI) at (TileH.30)  {\TileI};
			\node[hexsty,anchor=150] (TileJ) at (TileI.30)  {\TileJ};
			\node[hexsty,anchor=150] (TileK) at (TileJ.30)  {\TileK};
			\node[hexsty,anchor=150] (TileL) at (TileK.30)  {\TileL};
			
			\node[hexsty,anchor=90]  (TileM) at (TileH.330) {\TileM};
			\node[hexsty,anchor=150] (TileN) at (TileM.30)  {\TileN};
			\node[hexsty,anchor=150] (TileO) at (TileN.30)  {\TileO};
			\node[hexsty,anchor=150] (TileP) at (TileO.30)  {\TileP};
			
			\node[hexsty,anchor=90]  (TileQ) at (TileM.330) {\TileQ};
			\node[hexsty,anchor=150] (TileR) at (TileQ.30)  {\TileR};
			\node[hexsty,anchor=150] (TileS) at (TileR.30)  {\TileS};
			
			\node[portsty,anchor=300] (PortA) at (TileA.120) {\PortA};
			\node[portsty,anchor=240] (PortB) at (TileB.60)  {\PortB};
			\node[portsty,anchor=240] (PortC) at (TileG.60)  {\PortC};
			\node[portsty,anchor=180] (PortD) at (TileL.0)   {\PortD};
			\node[portsty,anchor=120] (PortE) at (TileP.300) {\PortE};
			\node[portsty,anchor=120] (PortF) at (TileR.300) {\PortF};
			\node[portsty,anchor=60]  (PortG) at (TileQ.240) {\PortG};
			\node[portsty,anchor=0]   (PortH) at (TileM.180) {\PortH};
			\node[portsty,anchor=0]   (PortI) at (TileD.180) {\PortI};
			
			\end{tikzpicture}
			
			 \StrSubstitute{#1}{.csv}{}{} % file name
			
		\end{minipage}%
		\hspace{-1.5cm}
		\begin{minipage}{.25\textwidth}
			\flushleft\noindent
			\begin{adjustbox}{angle=90}
				\begin{tabular}{@{}rr@{}}
					roads & 15\\
					sett.s &5\\
					cities& { }4\\
					& \\
					lgst road  & $\geq$ 5 \\
					& \\  % space under here to put dots with marker for lgst road
					lgst army  & $\geq$ 3\\
				\end{tabular}
			\end{adjustbox}
			
			\vspace{7cm}

			\begin{adjustbox}{angle=90}
				\begin{tabular}{@{}cll@{}}
					&& \\
					&& \\
					\RoadSym& road & \Port[n]{w} \Port[n]{b} \ExtraHeight\\
					\SettSym& sett.&\Port[n]{w} \Port[n]{b} \Port[n]{s} \Port[n]{h} \ExtraHeight\\
					\CitySym& city & \Port[n]{h} \Port[n]{h} \Port[n]{o} \Port[n]{o} \Port[n]{o} \ExtraHeight\\
					\DevSym &dev. & \Port[n]{s} \Port[n]{h} \Port[n]{o} \ExtraHeight\\
				\end{tabular}
			\end{adjustbox}
		\end{minipage}%
	}
	
	\tt% typewriter font
	\noindent%

	\renewcommand{\do}[1]{% this is performed in the csv do-loop	
		% load up commands that are pushed to the catan board tiles
		\edef\column{\xintthe\xintiexpr\value{index}/:3\relax}% get the column num of the csv

		\ifnum\column=0 % if first column, get the tile letter as a number
			\setcounter{tilenum}{\xintthe\xintiexpr(\value{index})/3+1\relax}%
		\fi%
		\ifnum\column=1 % second column is the resource, store it in a definition
			\def\resource{#1}%
		\fi%
		\ifnum\column=2 % third column is the number, now we can replace the Tile or Port Command
			\ifnum\value{index}<57
			
%				\Alph{tilenum}: \resource, #1\\
				\defineTileCmd{Tile\Alph{tilenum}}{#1}{\resource}%
			\else
			
				\addtocounter{tilenum}{-19}%
%				P\alph{tilenum}: \resource\\
				\definePortCmd{Port\Alph{tilenum}}{\resource}%
			\fi%
		\fi%
		\addtocounter{index}{1}% increment index
	}% new line in between items
	\newcommand{\defineTileCmd}[3]{%
	\expandafter\defineTileCmdAux\expandafter{#3}{#1}{#2}% expands the deine command and the \resource
	}
	\newcommand{\defineTileCmdAux}[3]{%
		\expandafter\newcommand\csname #2\endcsname{\Tile{#1}{#3}{}}%
	}
	\newcommand{\definePortCmd}[2]{%
	\expandafter\definePortCmdAux\expandafter{#2}{#1}%
	}
	\newcommand{\definePortCmdAux}[2]{%
		\expandafter\newcommand\csname #2\endcsname{\Port{#1}}%
	}


\foreach \file in \CSVfiles{% iterate through the CSV files
 	\CatchFileDef{\CSVdata}{\file}{\endlinechar=`\,}% ensure comma between lines
	\setcounter{index}{0}%  index is location in csv file
	\setcounter{column}{0}% 
	\setcounter{tilenum}{0}% number corresponding to tile letter (A=1, B=2, etc.)
	\expandafter\docsvlist\expandafter{\CSVdata} % run \do for CSV items, this will load commands to produce Catan Board
	\PrintCatanBoard[\file]% Print the board
}%


\end{document} % thanks :)